\documentclass[10pt,a4paper]{amsart}
\usepackage{amsmath,amssymb,amsthm,graphicx}
\usepackage{verbatim}
\newenvironment{reference}{\comment}{\endcomment}
%\newenvironment{reference}{}{}
\newenvironment{slogan}{\comment}{\endcomment}
\newenvironment{history}{\comment}{\endcomment}

% For commutative diagrams we use Xy-pic
\usepackage[all]{xy}

% We use 2cell for 2-commutative diagrams.
\xyoption{2cell}
% Theorem environments.
%
\theoremstyle{plain}
\newtheorem{theorem}[subsection]{Theorem}
\newtheorem{proposition}[subsection]{Proposition}
\newtheorem{lemma}[subsection]{Lemma}

\theoremstyle{definition}
\newtheorem{definition}[subsection]{Definition}
\newtheorem{example}[subsection]{Example}
\newtheorem{exercise}[subsection]{Exercise}
\newtheorem{situation}[subsection]{Situation}

\theoremstyle{remark}
\newtheorem{remark}[subsection]{Remark}
\newtheorem{remarks}[subsection]{Remarks}

\numberwithin{equation}{subsection}

% Macros
%
\def\lim{\mathop{\mathrm{lim}}\nolimits}
\def\colim{\mathop{\mathrm{colim}}\nolimits}
\def\Spec{\mathop{\mathrm{Spec}}}
\def\Hom{\mathop{\mathrm{Hom}}\nolimits}
\def\Ext{\mathop{\mathrm{Ext}}\nolimits}
\def\SheafHom{\mathop{\mathcal{H}\!\mathit{om}}\nolimits}
\def\SheafExt{\mathop{\mathcal{E}\!\mathit{xt}}\nolimits}
\def\Sch{\mathit{Sch}}
\def\Mor{\mathop{\mathrm{Mor}}\nolimits}
\def\Ob{\mathop{\mathrm{Ob}}\nolimits}
\def\Sh{\mathop{\mathit{Sh}}\nolimits}
\def\NL{\mathop{N\!L}\nolimits}
\def\CH{\mathop{\mathrm{CH}}\nolimits}
\def\proetale{{pro\text{-}\acute{e}tale}}
\def\etale{{\acute{e}tale}}
\def\QCoh{\mathit{QCoh}}
\def\Ker{\mathop{\mathrm{Ker}}}
\def\Im{\mathop{\mathrm{Im}}}
\def\Coker{\mathop{\mathrm{Coker}}}
\def\Coim{\mathop{\mathrm{Coim}}}

% Boxtimes
%
\DeclareMathSymbol{\boxtimes}{\mathbin}{AMSa}{"02}

%
% Macros for moduli stacks/spaces
%
\def\QCohstack{\mathcal{QC}\!\mathit{oh}}
\def\Cohstack{\mathcal{C}\!\mathit{oh}}
\def\Spacesstack{\mathcal{S}\!\mathit{paces}}
\def\Quotfunctor{\mathrm{Quot}}
\def\Hilbfunctor{\mathrm{Hilb}}
\def\Curvesstack{\mathcal{C}\!\mathit{urves}}
\def\Polarizedstack{\mathcal{P}\!\mathit{olarized}}
\def\Complexesstack{\mathcal{C}\!\mathit{omplexes}}
% \Pic is the operator that assigns to X its picard group, usage \Pic(X)
% \Picardstack_{X/B} denotes the Picard stack of X over B
% \Picardfunctor_{X/B} denotes the Picard functor of X over B
\def\Pic{\mathop{\mathrm{Pic}}\nolimits}
\def\Picardstack{\mathcal{P}\!\mathit{ic}}
\def\Picardfunctor{\mathrm{Pic}}
\def\Deformationcategory{\mathcal{D}\!\mathit{ef}}
\usepackage{cleveref}
\crefname{originallemma}{original lemma}{original lemmas}
\Crefname{originallemma}{Original lemma}{Original lemmas}
\makeatletter
\newlabel{algebra-lemma-ff-rings}{{00HQ}{source}{}{sourcepointer-1}{}}
\newlabel{algebra-lemma-ff-rings@cref}{{[originallemma][0][]00HQ}{[1][0][]source}{}{}{}}
\makeatother
\crefname{originalreference}{original reference}{original references}
\Crefname{originalreference}{Original reference}{Original references}
\makeatletter
\newlabel{definition-w-local}{{096A}{source}{}{sourcepointer-2}{}}
\newlabel{definition-w-local@cref}{{[originalreference][0][]096A}{[1][0][]source}{}{}{}}
\makeatother
\title{The w-localization of an affine scheme is ind-Zariski, faithfully flat and has the expected closed points and specializations}
\author{The Stacks Project Authors}
\date{stacks-project:proetale.tex:local-snapshot-20261012; GNU FDL 1.2 or later}
\begin{document}
\maketitle
\noindent\textit{Editorial note.} All delivered mathematical text is the author's own, copied from the pinned Stacks Project chapter file proetale.tex; no statement or proof step was rewritten, reordered or removed.\par\medskip
\noindent\textit{Editorial note.} The counted proof is the author's complete argument: the derivation of faithful flatness from surjectivity, the description of the closed points through the closures of the images in the partial w-localizations, the two-point separation argument proving uniqueness of specializations into Z, and the reducedness computation.\par\medskip
\noindent\textit{Editorial note.} The label lemma-make-w-local names the counted statement, and the proof environment immediately following it in the source is delivered verbatim as its proof.\par\medskip
\noindent\textit{Editorial note.} The counted statement is carried by the author's own lemma environment in the chapter source.\par\medskip
\noindent\textit{Editorial note.} Pre-open gate: python operations/agent-authoring/tools/precard\_gate.py --source proetale.tex --label lemma-make-w-local --start 843 --end 900 --local-id 3578 returned RESULT=OK (L3578 free, occupancy FREE, public interval clear). Candidate taken from sources/stacks-proetale/candidates-halfA.json.\par\medskip
\noindent\textit{Editorial note.} The proof refers to 'the reader verifies' for the elementary nilpotent computation that finishes part (4); that computation is a one-line verification of the displayed description of the Z\_E rather than an unproved mathematical step, so the statement is counted with the phrase recorded here verbatim.\par\medskip
\noindent\textit{Editorial note.} The statement and proof are copied from the pinned chapter file; the cited inputs are kept as external pointers with their canonical Stacks tags.\par\medskip
\noindent\textit{Allowed standard prerequisites (editorial).}\begin{itemize}
\item Faithful flatness from surjectivity of the spectrum map. Algebra, the lemma that a ring map is faithfully flat if the induced map on spectra is surjective and the map is flat. Supplies the faithful flatness of A -> A\_w once surjectivity on spectra is known.
\item Structure of the partial w-localizations Z\_E. Same chapter, the construction of the Z\_E as disjoint unions of locally closed subschemes isomorphic to spectra of localizations of quotients. Supplies the description of Z\_E used in the proof that the Z\_E have no nilpotents and in the bijectivity Z -> X.
\item Quasi-finiteness and open images for etale maps. Algebra, the lemmas on the local structure of the partial w-localizations and the finiteness of the primes lying over a maximal ideal. Supplies the identification of the closed points of the limit and the absence of specializations between the points of Z.
\end{itemize}\medskip
\ifdefined\hypertarget\hypertarget{sourcepointer-1}{}\fi
\noindent\textit{Source reference.} Original lemma 00HQ: sources/stacks-chow/tags.txt, source lines 652--652. Prerequisite pointer invoked at proof line(s) 862; it names a standard fact of a sibling chapter that the counted argument uses, and every step of the counted proof built on it is retained verbatim in the delivered body.\par\medskip
\ifdefined\hypertarget\hypertarget{sourcepointer-2}{}\fi
\noindent\textit{Source reference.} Original reference 096A: sources/stacks-proetale/pinned-original-source.tex, source lines 405--405. Prerequisite pointer invoked at proof line(s) 854; it names a standard fact of the same chapter that the counted argument uses, and every step of the counted proof built on it is retained verbatim in the delivered body.\par\medskip
\begin{lemma}
\label{lemma-make-w-local}
Let $X = \Spec(A)$ be an affine scheme. With $A \to A_w$, $X_w = \Spec(A_w)$,
and $Z \subset X_w$ as above.
\begin{enumerate}
\item $A \to A_w$ is ind-Zariski and faithfully flat,
\item $X_w \to X$ induces a bijection $Z \to X$,
\item $Z$ is the set of closed points of $X_w$,
\item $Z$ is a reduced scheme, and
\item every point of $X_w$ specializes to a unique point of $Z$.
\end{enumerate}
In particular, $X_w$ is w-local (Definition \ref{definition-w-local}).
\end{lemma}

\begin{proof}
The map $A \to A_w$ is ind-Zariski by construction.
For every $E$ the morphism $Z_E \to X$ is a bijection, hence (2).
As $Z \subset X_w$ we conclude $X_w \to X$ is surjective and
$A \to A_w$ is faithfully flat by
Algebra, Lemma \ref{algebra-lemma-ff-rings}. This proves (1).

\medskip\noindent
Suppose that $y \in X_w$, $y \not \in Z$. Then there
exists an $E$ such that the image of $y$ in $X_E$ is not contained in
$Z_E$. Then for all $E \subset E'$ also $y$ maps to an element of $X_{E'}$
not contained in $Z_{E'}$. Let $T_{E'} \subset X_{E'}$ be the reduced
closed subscheme which is the closure of the image of $y$. It is
clear that $T = \lim_{E \subset E'} T_{E'}$ is the closure of $y$ in $X_w$.
For every $E \subset E'$ the scheme $T_{E'} \cap Z_{E'}$ is nonempty
by construction of $X_{E'}$. Hence $\lim T_{E'} \cap Z_{E'}$ is nonempty
and we conclude that $T \cap Z$ is nonempty. Thus $y$ is not a closed point.
It follows that every closed point of $X_w$ is in $Z$.

\medskip\noindent
Suppose that $y \in X_w$ specializes to $z, z' \in Z$. We will show that
$z = z'$ which will finish the proof of (3) and will imply (5).
Let $x, x' \in X$ be the images of $z$ and $z'$. Since $Z \to X$ is
bijective it suffices to show that $x = x'$. If $x \not = x'$, then
there exists an $f \in A$ such that $x \in D(f)$ and $x' \in V(f)$
(or vice versa). Set $E = \{f\}$ so that
$$
X_E = \Spec(A_f) \amalg \Spec(A_{V(f)}^\sim)
$$
Then we see that $z$ and $z'$ map $x_E$ and $x'_E$ which are in different
parts of the given decomposition of $X_E$ above. But then it impossible
for $x_E$ and $x'_E$ to be specializations of a common point.
This is the desired contradiction.

\medskip\noindent
Recall that given a finite subset $E \subset A$ we have $Z_E$
is a disjoint union of the locally closed subschemes $Z(E', E'')$
each isomorphic to the spectrum of $(A/I)_f$ where $I$ is the ideal
generated by $E''$ and $f$ the product of the elements of $E'$.
Any nilpotent element $b$ of $(A/I)_f$ is the class of $g/f^n$
for some $g \in A$. Then setting $E' = E \cup \{g\}$ the reader
verifies that $b$ is pulls back to zero under the transition map
$Z_{E'} \to Z_E$ of the system. This proves (4).
\end{proof}

\end{document}
